\section{The Words We Shall Answer For}

Christ expressly includes our speech in the judgment. After speaking of the tree known by its fruit and the mouth speaking from the heart's abundance, he warns that we shall render an account for every idle word (Mt 12:33--37). Words are acts by which the inward life enters the world. Through them we acknowledge God, instruct a child, bind ourselves by a promise, comfort the afflicted, accuse the innocent, or persuade another person to sin. The tongue's work can pass quickly and leave consequences which last for years.

The warning occurs after Christ has delivered a man from demonic affliction and his opponents have attributed the work to the prince of demons. Their charge interprets a manifest good through a will set against its author. Chrysostom attends closely to this contradiction (\emph{Homilies on Matthew} XLII.1). They do not deny that the sufferer has been healed, yet they malign the healer. The scene exposes a form of falsehood more deliberate than an error in information: a refusal to let another person's good work disturb a settled hostility.

We can recognize a smaller version of that refusal in our own speech. Someone we dislike acts generously, and we immediately supply a selfish motive. A person we resent tells the truth, and we find a way to make the truth sound contemptible. We may possess no evidence for the explanation, but it preserves the judgment we prefer. Such habits teach the heart to resist gratitude. Before Christ, the question will include what we did with the good we saw, and whether envy made us unwilling to acknowledge it.

Chrysostom explains the idle word here through falsehood and unjust accusation; he also records the interpretation which includes empty, coarse, or immodest speech. A person must answer for what he himself has said (XLII.1). A passing remark can contain a deliberate injustice. We may repeat an accusation without caring whether it is true, or exaggerate a failure to obtain laughter at someone else's expense. Examination must consider the intention with which we spoke as well as the words we used.

James describes the disproportion between the tongue's size and its power through a bit, a rudder, and a little fire which sets a great wood ablaze (Jas 3:1--6). A few words can direct a household's judgment of an absent relative or establish the way a community speaks about a vulnerable person. When the story spreads, its first speaker may lose control of it. Even after he has forgotten the remark, others may repeat it as a reason for despising someone whom God has commanded them to love.

This power also belongs to written speech. A letter, a message, or a public comment can preserve the moment in which anger outran truth and charity. Ease of transmission does not remove the speaker from the act. We can send an accusation to many people before taking the time to learn whether it is true, then treat the resulting injury as something caused by the audience. A sober preparation for judgment includes accepting authorship of our own words and the foreseeable harm we willingly set in motion.

James then turns from the tongue's power to its contradiction. With it we bless God and curse human beings made in his likeness (Jas 3:9--12). His objection concerns the relation of both acts to God. In worship we acknowledge the Creator; in cursing our neighbor we dishonor someone created in his image. The obligation to reverence that image remains when prayer has ended. A conversation afterward is one place to fulfill it, especially when someone begins speaking contemptuously of an absent person.

The Church gives precise names to several ways in which speech violates this obedience. Rash judgment assumes another's moral fault without sufficient foundation; detraction needlessly exposes real faults to people unaware of them; calumny injures reputation through falsehood (CCC 2477). These distinctions help us examine ourselves honestly. Saying that a story is true does not by itself justify passing it on. Saying that we merely repeated what we heard does not remove the responsibility to care whether our words injure someone unjustly. Truth and charity concern both the content and the purpose of communication.

There are occasions when truth must be spoken to protect someone or fulfill a duty of justice. A witness may need to report wrongdoing, a parent to confront a danger, or a person entrusted with care to seek help from someone able to intervene. Such speech has an end which idle accusation lacks: the good of persons and the right ordering of a real responsibility. The Catechism requires charity and respect for truth in deciding what should be communicated (CCC 2488--2489). Prudence seeks the appropriate person, manner, and extent; it does not turn concern for reputation into permission to conceal a continuing injury.

The temptation to silence can therefore be as searching as the temptation to gossip. We may dislike conflict, fear losing approval, or benefit from leaving an injustice undisturbed. If an actual duty requires us to speak, a pleasant appearance of peace can conceal a refusal of courage. Christ's judgment includes omissions of truthful witness as well as malicious words. Preparation asks us to submit speech and silence alike to charity, accepting that the right word may cost something and that the satisfying word may need to remain unspoken.

Chrysostom gives a particular consolation to the slandered. They will not have to answer for the evil another person falsely says about them; the slanderer will answer for the wrong he has committed (XLII.2). The injured person may still need to defend himself, and the injury can be painful even when the accusation is plainly false. But he can pursue a correction without treating the false report as a verdict from God. His own account concerns what he freely did, including his response to the injury.

For the person who has harmed another by speech, repentance reaches the damage done. Forgiveness leaves a duty of reparation where justice requires it (CCC 2487). We may need to correct a false statement before those who heard it, withdraw an unfair judgment, or acknowledge that we spoke without knowing. The necessary repair may be embarrassing precisely because the original remark enhanced our standing. That cost belongs to the return to truth. We cannot always undo the whole effect, but we can cease lending our authority to the wrong and seek sound guidance about what further repair is possible.

Good speech also enters the account of a life. A patient explanation can spare someone a lasting misunderstanding. When a parent acknowledges a fault and repairs it, a child can learn how repentance leads to amendment. Faithful testimony to Christ may help someone long after the speaker has forgotten giving it. Christ's account of good fruit from a good heart gives us reason to cultivate such speech. Restraint leaves us freer to say what another person actually needs to hear.

Chrysostom follows his treatment of harmful speech by directing attention back to personal examination. He imagines conscience being summoned to consider its own acts instead of escaping into accusations against others (XLII.3). We can spend an evening rehearsing another person's failure and overlook the anger or vanity with which we rehearsed it. Examination asks about our own choice: what we said, what we knew, and what we intended. These matters can be brought to confession while amendment remains possible.

The answer will sometimes require more than a resolution to speak carefully. If resentment continually supplies our conversation, the resentment itself needs conversion. If we must always appear knowledgeable, the vanity which makes us pretend knowledge needs surrender. If contempt wins us a place among friends, we may have to accept the loss of that approval. Christ's image of the tree directs attention to the root. Grace can make the heart more truthful and charitable, so that restraint grows into a different manner of speaking.

Annual recollection can uncover a habit which a single day's examination misses. We may repeatedly return to one humiliating story or speak contemptuously of someone whenever a certain group gathers. That recurrence may require a deliberate change in the company we keep or the conversations we permit ourselves. Words of gratitude and reconciliation may also have been delayed for months. Where a duty to speak has become clear, we can decide when and how to fulfill it.

